% GENERATED FILE. Edit canonical lessons, then run npm run generate:books.
% canonical-source-hash: fnv1a64:7d05ecbba9d79cce
% canonical-lessons: UR-C189-jab-tak, UR-C189-tab-tak, UR-C189-abhi-tak, UR-C189-isi-vaqt, UR-C189-der-se

\chapter{Until, Till then, So far, Right away, Late}
\label{ch:ur-until-till-then-so-far-right-away-late}
\hlchaptermodality{189}

\begin{chapteropening}
\textbf{By the end of this chapter:} \emph{I can say until, till then, so far, right away and late.}

Complete the last lesson of chapter 189: I can say until, till then, so far, right away and late.
\end{chapteropening}

\section[\texorpdfstring{\ur{جب} \ur{تک} (jab tak)}{جب تک (jab tak)}]{\ur{جب} \ur{تک} (jab tak) — until; as long as}

\label{lesson:UR-C189-jab-tak}



\begin{quote}
Before the new one: say the Urdu for July, then the Urdu for August.
\end{quote}

\subsection*{You'll want to know: \ur{جب} \ur{تک}}
\textbf{\ur{جب} \ur{تک}} — \emph{jab tak} — \textquotedblleft{}until; as long as\textquotedblright{}.

\begin{cousinweb}
\textbf{\ur{جب} \ur{تک}} (\emph{jab tak}) is \textquotedblleft{}when, up to\textquotedblright{}. CDIAL lists \emph{jab}, \textquotedblleft{}when\textquotedblright{}, under Sanskrit \emph{evam eva}, \textquotedblleft{}exactly so\textquotedblright{}, crossed with the relative stem \emph{ya-} (CDIAL 2528), and \emph{tak}, \textquotedblleft{}till\textquotedblright{}, under Sanskrit \emph{tāvatka}, built on \emph{tāvant}, \textquotedblleft{}so much\textquotedblright{} (CDIAL 5805).
\end{cousinweb}

\begin{grammarlens}[title={What you've built}]
One more word for this chapter.
\end{grammarlens}

\subsection*{Your turn}
\begin{itemize}
\raggedright
  \item \emph{Say it:} \emph{jab tak}
  \item \emph{Say it:} \emph{jab tak}, once more
  \item \emph{Say it:} say \emph{jab tak}
  \item \emph{Recall:} say the Urdu for July, then the Urdu for August, then say \emph{jab tak} again
\end{itemize}

\begin{tcolorbox}[breakable,colback=teal!4,colframe=teal!35!black,title={Before you move on}]
What does \ur{جب} \ur{تک} mean? (\textquotedblleft{}Until; as long as\textquotedblright{}.) Say it once more.
\end{tcolorbox}
\section[\texorpdfstring{\ur{تب} \ur{تک} (tab tak)}{تب تک (tab tak)}]{\ur{تب} \ur{تک} (tab tak) — till then}

\label{lesson:UR-C189-tab-tak}



\begin{quote}
Before the new one: say the Urdu for August, then the Urdu for until; as long as.
\end{quote}

\subsection*{You'll want to know: \ur{تب} \ur{تک}}
\textbf{\ur{تب} \ur{تک}} — \emph{tab tak} — \textquotedblleft{}till then\textquotedblright{}.

\begin{cousinweb}
\textbf{\ur{تب} \ur{تک}} (\emph{tab tak}) is \textquotedblleft{}then, up to\textquotedblright{}. CDIAL lists \emph{tab}, \textquotedblleft{}then\textquotedblright{}, under Sanskrit \emph{evam eva}, \textquotedblleft{}exactly so\textquotedblright{}, crossed with the pointing stem \emph{ta-} (CDIAL 2528), and \emph{tak}, \textquotedblleft{}till\textquotedblright{}, under Sanskrit \emph{tāvatka}, built on \emph{tāvant}, \textquotedblleft{}so much\textquotedblright{} (CDIAL 5805).
\end{cousinweb}

\begin{grammarlens}[title={What you've built}]
One more word for this chapter.
\end{grammarlens}

\subsection*{Your turn}
\begin{itemize}
\raggedright
  \item \emph{Say it:} \emph{tab tak}
  \item \emph{Say it:} \emph{tab tak}, once more
  \item \emph{Say it:} say \emph{tab tak}
  \item \emph{Recall:} say the Urdu for August, then the Urdu for until; as long as, then say \emph{tab tak} again
\end{itemize}

\begin{tcolorbox}[breakable,colback=teal!4,colframe=teal!35!black,title={Before you move on}]
What does \ur{تب} \ur{تک} mean? (\textquotedblleft{}Till then\textquotedblright{}.) Say it once more.
\end{tcolorbox}
\section[\texorpdfstring{\ur{ابھی} \ur{تک} (abhī tak)}{ابھی تک (abhī tak)}]{\ur{ابھی} \ur{تک} (abhī tak) — so far, still}

\label{lesson:UR-C189-abhi-tak}



\begin{quote}
Before the new one: say the Urdu for until; as long as, then the Urdu for till then.
\end{quote}

\subsection*{You'll want to know: \ur{ابھی} \ur{تک}}
\textbf{\ur{ابھی} \ur{تک}} — \emph{abhī tak} — \textquotedblleft{}so far, still\textquotedblright{}.

\begin{cousinweb}
\textbf{\ur{ابھی} \ur{تک}} (\emph{abhī tak}) is \textquotedblleft{}right now, up to\textquotedblright{}. CDIAL lists \emph{abhī} under Sanskrit \emph{evam eva}, \textquotedblleft{}exactly so\textquotedblright{}, crossed with the pointing stem \emph{a-} (CDIAL 2528), and \emph{tak}, \textquotedblleft{}till\textquotedblright{}, under Sanskrit \emph{tāvatka}, built on \emph{tāvant}, \textquotedblleft{}so much\textquotedblright{} (CDIAL 5805).
\end{cousinweb}

\begin{grammarlens}[title={What you've built}]
One more word for this chapter.
\end{grammarlens}

\subsection*{Your turn}
\begin{itemize}
\raggedright
  \item \emph{Say it:} \emph{abhī tak}
  \item \emph{Say it:} \emph{abhī tak}, once more
  \item \emph{Say it:} say \emph{abhī tak}
  \item \emph{Recall:} say the Urdu for until; as long as, then the Urdu for till then, then say \emph{abhī tak} again
\end{itemize}

\begin{tcolorbox}[breakable,colback=teal!4,colframe=teal!35!black,title={Before you move on}]
What does \ur{ابھی} \ur{تک} mean? (\textquotedblleft{}So far, still\textquotedblright{}.) Say it once more.
\end{tcolorbox}
\section[\texorpdfstring{\ur{اسی} \ur{وقت} (isī vaqt)}{اسی وقت (isī vaqt)}]{\ur{اسی} \ur{وقت} (isī vaqt) — right away, immediately}

\label{lesson:UR-C189-isi-vaqt}



\begin{quote}
Before the new one: say the Urdu for till then, then the Urdu for so far.
\end{quote}

\subsection*{You'll want to know: \ur{اسی} \ur{وقت}}
\textbf{\ur{اسی} \ur{وقت}} — \emph{isī vaqt} — \textquotedblleft{}right away, immediately\textquotedblright{}.

\begin{cousinweb}
\textbf{\ur{اسی} \ur{وقت}} (\emph{isī vaqt}) is \textquotedblleft{}at this very time\textquotedblright{}. \emph{isī} is \emph{is}, \textquotedblleft{}this\textquotedblright{}, with the emphatic \emph{-ī}; CDIAL lists \emph{is} under Sanskrit \emph{eṣa}, \textquotedblleft{}this\textquotedblright{} (CDIAL 2530). \emph{vaqt} is Arabic \emph{waqt}, \textquotedblleft{}a time\textquotedblright{}, which reached Urdu through Persian (Lane, root \emph{w-q-t}).
\end{cousinweb}

\begin{grammarlens}[title={What you've built}]
One more word for this chapter.
\end{grammarlens}

\subsection*{Your turn}
\begin{itemize}
\raggedright
  \item \emph{Say it:} \emph{isī vaqt}
  \item \emph{Say it:} \emph{isī vaqt}, once more
  \item \emph{Say it:} say \emph{isī vaqt}
  \item \emph{Recall:} say the Urdu for till then, then the Urdu for so far, then say \emph{isī vaqt} again
\end{itemize}

\begin{tcolorbox}[breakable,colback=teal!4,colframe=teal!35!black,title={Before you move on}]
What does \ur{اسی} \ur{وقت} mean? (\textquotedblleft{}Right away, immediately\textquotedblright{}.) Say it once more.
\end{tcolorbox}
\section[\texorpdfstring{\ur{دیر} \ur{سے} (der se)}{دیر سے (der se)}]{\ur{دیر} \ur{سے} (der se) — late}

\label{lesson:UR-C189-der-se}



\begin{quote}
Before the new one: say the Urdu for so far, then the Urdu for right away.
\end{quote}

\subsection*{You'll want to know: \ur{دیر} \ur{سے}}
\textbf{\ur{دیر} \ur{سے}} — \emph{der se} — \textquotedblleft{}late\textquotedblright{}.

\begin{cousinweb}
\textbf{\ur{دیر} \ur{سے}} (\emph{der se}) is \textquotedblleft{}from lateness\textquotedblright{}. \emph{der} is usually traced to Persian \emph{dēr}, \textquotedblleft{}late\textquotedblright{}. CDIAL lists \emph{se} under Sanskrit \emph{sahita}, \textquotedblleft{}together with\textquotedblright{} (CDIAL 13310), and \emph{sama}, \textquotedblleft{}equal\textquotedblright{}, has also been suggested (CDIAL 13173).
\end{cousinweb}

\begin{grammarlens}[title={What you've built}]
One more word for this chapter.
\end{grammarlens}

\subsection*{Your turn}
\begin{itemize}
\raggedright
  \item \emph{Say it:} \emph{der se}
  \item \emph{Say it:} \emph{der se}, once more
  \item \emph{Say it:} say \emph{der se}
  \item \emph{Recall:} say the Urdu for so far, then the Urdu for right away, then say \emph{der se} again
\end{itemize}

\begin{tcolorbox}[breakable,colback=teal!4,colframe=teal!35!black,title={Before you move on}]
What does \ur{دیر} \ur{سے} mean? (\textquotedblleft{}Late\textquotedblright{}.) Say it once more.
\end{tcolorbox}
